\subsection*{Summary and interpretation.}
\noindent
At a high level, our data show that brief, low-intensity tFUS to the sgACC increases its functional coupling with the rest of the brain in the period \emph{after} sonication, relative to a sham control (e.g., Pre→Post Active–Sham $\Delta$: $\beta=0.079$, 95\% CI [0.006, 0.151], $p=0.033$). This effect emerged in a conservative, subject-level mixed-effects model that (i) respects the repeated-measures design, (ii) avoids edge-wise pseudoreplication, and (iii) does not rely on equal baselines, instead estimating a difference-in-differences (DiD) effect between Active and Sham sessions. We did not detect a comparably strong differential effect \emph{during} sonication, suggesting a delayed or accumulating influence of stimulation rather than an acute ``online'' change.

Analysis of the sgACC's connectivity with canonical brain networks provides insight into the nature of this accumulating effect. Over the course of the Sham session, sgACC's connectivity with the DMN increased over time -- consistent with the tendency for unconstrained rest to drift toward mind-wandering and internally oriented cognition \cite{Raichle2001,Mason2007,Christoff2009}. In contrast, Active tFUS did not enhance the sgACC's coupling with the DMN; instead, its connectivity increases with \emph{non-DMN} systems, most prominently the frontoparietal Cognitive Control network and, to a lesser extent, Salience and Somatomotor networks. The Cognitive Control network is established as a flexible hub supporting goal-directed regulation and cognitive reappraisal \cite{yeo2011organization,cole2013multi}. The disassociation found here -- increasing DMN coupling absent stimulation versus enhanced integration with non-DMN networks during Active sessions —- argues that tFUS is not bluntly amplifying the sonicated region's activity. Rather, it appears to bias sgACC activity toward circuits implicated in top-down regulation and away from those supporting self-referential processing. Relative to Sham, sgACC-DMN connectivity actually decreased over time, which is interesting in light of the evidence that sgACC-DMN hyperconnectivity tracks rumination and symptom burden in depression \cite{whitfield2012default}.

A dynamic systems interpretation of our findings is that tFUS biases ongoing dynamics toward distinct regimes. Resting-state activity naturally explores a landscape in which DMN-dominated states are common; a small, focal perturbation may inject sufficient exogenous drive to nudge the system toward alternative basins of attraction that prioritize control and monitoring. Although we did not measure metabolism, this account harmonizes qualitatively with longstanding observations of sgACC metabolic abnormalities in mood disorders and their modulation with successful intervention \cite{Mayberg2005,Drevets2008}. If tFUS can transiently reweight sgACC interactions toward control-related circuits, it may counteract the spontaneous drift into DMN-dominant modes that accompany passive rest. This is necessarily a mechanistic hypothesis (i.e., not a conclusion) but it generates concrete predictions for future multimodal work (e.g., Magnetic Resonance Spectroscopy or Positron Emission Tomography alongside tFUS). 

A critical aspect of our findings is the enhancement of FC \emph{following} stimulation. This aligns with prior reports of prolonged tFUS effects in humans and nonhuman primates \cite{Verhagen2019-ae,zhang2023excitatory,yaakub2023transcranial} and is consistent with processes that outlast the sonication epoch. Multiple biophysical pathways have been proposed for ultrasound–brain interactions, including modulation of mechanosensitive ion channels \cite{kubanek2016ultrasound}, changes to effective membrane mechanics \cite{krasovitski2011intramembrane}, and astrocytic contributions \cite{yang2015enhancement}. The present BOLD–FC measures cannot adjudicate among these mechanisms. The slow, post-sonication increase observed here suggests a delayed or accumulating influence of stimulation on network organization. We therefore view the persistence of the effect as hypothesis-generating evidence for longer-timescale neuromodulatory or plasticity-consistent processes. 

In order to rule out the possibility that baseline connectivity differences between conditions spuriously led to the observed post-sonication connectivity enhancement, we formally analyzed the relationship between the ``starting point'' of each session and the subsequent change in connectivity.  Under a pure regression-to-the-mean account, both conditions should exhibit a similar inverse relationship between baseline FC and the subsequent change. Instead, we observe a robust negative slope in Sham but a mildly \emph{positive} slope in Active, with their difference exhibiting statistical significance (Sham: slope $=-0.78$, CI $[-1.37, -0.18]$; Active: slope $=+0.21$, CI $[-0.58, 0.99]$). This is inconsistent with a regression-to-the-mean explanation, and indicates that  active tFUS alters the baseline--change relationship itself. 

Modeling scalp-to-sgACC distance as a continuous moderator yielded a directionally consistent, but trend-level, attenuation of the connectivity enhancement with increasing distance ($p=0.083$). The local regression analysis hints at an inverted-U shape between distance and effect size (Fig \ref{fig:distance_loess}). This pattern accords with a dose-attenuation hypothesis, where subjects whose sgACC is closest to the beam focus experience the largest increase in FC. Nevertheless, in our $n=16$ sample cohort, the post-stimulation effect persists after adjusting for distance, indicating that distance alone does not account for the observed FC changes. Future work incorporating subject-specific acoustic modeling (skull thickness and density), improved targeting uncertainty estimates, and direct intensity proxies may sharpen the sensitivity of distance (or dose) as a moderator.


\subsection*{Limitations}
\noindent
Our study is not without limitations. Primarily, our modest sample size likely limited our ability to resolve significant effects during sonication, as well as sgACC coupling increases with non-DMN networks other than the Cognitive Control network. In particular, the sgACC-Salience and sgACC-Somatomotor connectivity increase (Table \ref{tab:network_specificity_onesided}) did not survive FDR but showed uncorrected trends (e.g., Salience q=0.08; Somatomotor q=0.07). Secondly, the BOLD signal is challenging to interpret due to the complex biophysics that give rise to its generation \cite{buxton2013physics}. In particular, neurovascular coupling is largely agnostic to excitatory versus inhibitory neurotransmission \cite{logothetis2008we}, meaning that it is difficult to infer the effect of tFUS on cortical excitability and the overall direction of the change in neural activity. Our employment of FC as the dependent variable -- motivated by its well-established association with brain state -- further exacerbates this problem, as an increase in connectivity does not imply increased activation of the regions comprising the FC measurement. Alternative forms of MRI such as Magnetic Resonance Spectroscopy (MRS) \cite{yaakub2023transcranial} could permit a clearer understanding of the relative effects of tFUS on excitation and inhibition. Similarly, the utilization of concurrent electrophysiological recordings \cite{kim2023transcranial,mueller2014transcranial,kim2022neuromodulation,lee2016transcranial,min2011focused} may afford more direct insight into how sonication modulates neural activity. Finally, our study employed a cohort of healthy adult subjects lacking a clinical diagnosis of major depression or other mood disorders associated with impaired sgACC function. As such, it is unclear whether the modulation of sgACC connectivity found here translates to brains whose activity is dysregulated by an affective disorder. Future studies investigating tFUS in clinical populations with aberrant sgACC dynamics are motivated by the present study. 



